\subsection{P006 --- Understanding the effect of outdoor pollution episodes and HVAC type on indoor air quality}
\label{paper:P006}
\textbf{Authors:} Tristalee Mangin, Zachary Barrett, Zachary Palmer, Dillon Tang, Sean Nielson, Darrah Sleeth, Kerry Kelly\par
\textbf{Major Category:} Building Environment and IEQ\par
\textbf{Secondary Category:} HVAC and Building Services, Sensors, IoT and Data Integration\par
\textbf{Keywords:} PM2.5, HVAC, Infiltration, Wildfire smoke, Indoor air quality, Low-cost sensors\par
\paragraph{主な主張}
18か月間に21台のlow-cost PM2.5 sensorを用い, inversion, dust, wildfire smoke時のoutdoor-to-indoor infiltrationとHVAC typeの影響を評価した. 平均infiltration factorは0.07, 0.10, 0.37で, air-side economizer有りのHVACは他方式より高いinfiltrationを示した. \cite{Mangin:2025OutdoorPM25}
\paragraph{新規性}
異なるoutdoor pollution episodeを同一のlong-term monitoringで比較し, PM2.5 infiltrationをHVAC system typeと結び付けて評価した点に特徴がある.
\paragraph{位置付け}
outdoor PM2.5の室内侵入とHVACの影響が主要対象であるためBuilding Environment and IEQを主分類とし, HVACおよびsensor monitoringを副分類とする.
